\chapter{Introdução}
\label{cap:introducao}

A crescente digitalização de processos sociais, científicos, industriais e organizacionais tem provocado um aumento contínuo na produção, armazenamento, transmissão e processamento de dados. Esse cenário intensifica a relevância dos sistemas intensivos em dados, isto é, sistemas projetados para lidar com grandes volumes de dados, múltiplas fontes de informação, diferentes formatos e requisitos rigorosos de desempenho, escalabilidade, disponibilidade e confiabilidade. Em ambientes de Big Data, esses desafios são frequentemente associados às características de volume, velocidade, variedade e veracidade dos dados, as quais influenciam diretamente o modo como as arquiteturas de software são concebidas e avaliadas \cite{dai2019big,hilbert2016big}.

Diferentemente de sistemas tradicionais, aplicações intensivas em dados dependem da integração entre componentes distribuídos, mecanismos de armazenamento, plataformas de processamento paralelo, sistemas de mensageria, estratégias de movimentação de dados e técnicas de gerenciamento de recursos. Nesse contexto, decisões arquiteturais deixam de ser apenas escolhas estruturais de alto nível e passam a exercer impacto direto sobre atributos de qualidade, como latência, \textit{throughput}, escalabilidade, tolerância a falhas e eficiência no uso de recursos computacionais \cite{mattmann2004software,ji2020quality,kleppmann2017designing}.

Entre as decisões arquiteturais mais relevantes nesse domínio, destacam-se aquelas relacionadas à forma como os dados são ingeridos, movimentados, armazenados e processados ao longo de um \textit{pipeline}. A escolha entre processamento em lote (\textit{batch processing}) e processamento em fluxo (\textit{stream processing}), por exemplo, pode alterar substancialmente o comportamento do sistema em termos de vazão, tempo de resposta, estabilidade temporal e consumo de recursos. Do mesmo modo, decisões sobre particionamento, escalonamento, localização do processamento, uso de memória, transferência entre sistemas e configuração de mecanismos de tolerância a falhas podem produzir impactos distintos conforme o volume dos dados, a taxa de chegada dos eventos e o ambiente de execução \cite{haynes2016pipegen,wang2018hard,aung2019coordinate,lyu2022fine}.

Apesar da existência de diversas técnicas, \textit{frameworks} e estratégias voltadas à otimização de sistemas intensivos em dados, a literatura ainda apresenta evidências empíricas fragmentadas. Muitos estudos avaliam soluções específicas em contextos particulares, utilizando métricas, \textit{workloads} e ambientes experimentais heterogêneos. Essa fragmentação dificulta a comparação entre abordagens, reduz a capacidade de generalização dos resultados e limita o apoio sistemático à tomada de decisão arquitetural. Como consequência, arquitetos de software frequentemente precisam decidir entre alternativas tecnológicas e arquiteturais com base em experiências prévias, recomendações práticas ou experimentos isolados, sem dispor de mecanismos quantitativos capazes de orientar essas escolhas de forma mais objetiva.

Diante desse contexto, esta pesquisa parte do seguinte problema: como apoiar, de forma sistemática e baseada em evidências, a tomada de decisão arquitetural em sistemas intensivos em dados, considerando características dos dados, estratégias de processamento e métricas observáveis de desempenho?

Para enfrentar essa lacuna, a pesquisa foi organizada em três movimentos complementares. O primeiro consistiu na realização de uma Revisão Sistemática da Literatura sobre o \textit{design} de sistemas intensivos em dados, acompanhada de uma síntese quantitativa exploratória. Essa etapa permitiu identificar técnicas arquiteturais recorrentes, métricas utilizadas na avaliação de desempenho, padrões de otimização e lacunas metodológicas relacionadas ao relato experimental.

O segundo movimento consistiu na condução de uma campanha experimental voltada à produção de evidência empírica controlada sobre decisões arquiteturais em sistemas intensivos em dados. O experimento comparou os paradigmas de processamento \textit{batch} e \textit{streaming} utilizando tecnologias amplamente adotadas, como Apache Spark, Spark Structured Streaming e Apache Kafka. Foram avaliados diferentes volumes de dados, taxas de eventos, números de \textit{workers} e configurações de \textit{trigger}, considerando métricas como \textit{throughput}, latência, uso de CPU, consumo de memória, \textit{backpressure} e taxa de processamento \cite{zaharia2016apache,armbrust2018structured,kreps2011kafka,matteussi2022performance}.

Os resultados experimentais obtidos reforçam a importância de compreender os \textit{trade-offs} envolvidos nas decisões arquiteturais. O processamento \textit{batch} tende a favorecer maiores taxas de \textit{throughput} em determinados cenários, enquanto o processamento em fluxo introduz custos associados à ingestão contínua, à coordenação dos micro-lotes e ao controle de estabilidade temporal. Além disso, fatores como volume de dados, paralelismo horizontal e intervalo de \textit{trigger} podem alterar significativamente o comportamento do sistema, indicando que decisões arquiteturais devem ser analisadas de forma contextualizada e baseada em evidências.

O terceiro movimento da pesquisa consiste na proposta de uma abordagem quantitativa de apoio à decisão arquitetural, considerando evidências empíricas e características informacionais dos dados. A intenção é investigar se propriedades como volume, taxa de eventos, entropia, redundância e variabilidade podem ser combinadas com métricas de desempenho para apoiar a comparação entre alternativas arquiteturais.


A Figura~\ref{fig:fluxo-pesquisa} apresenta o encadeamento geral da pesquisa.

\begin{figure}[htbp]
\centering
\begin{tikzpicture}[
    node distance=1.2cm and 1.2cm,
    etapa/.style={
        rectangle, rounded corners, draw, align=center,
        text width=3.2cm, minimum height=1.2cm,
        font=\small\sffamily, fill=gray!5
    },
    destaque/.style={
        etapa, font=\small\bfseries, fill=blue!10, text width=4.5cm
    },
    seta/.style={-{Latex[length=3mm]}, thick, draw=gray!80}
]

% Nós
\node[etapa] (rsl) {Revisão Sistemática da Literatura};
\node[etapa, right=of rsl] (exp) {Campanha Experimental};
\node[etapa, right=of exp] (modelo) {Modelo de Apoio à Decisão};
\node[destaque, below=1.5cm of exp] (decisao) {Tomada de Decisão Arquitetural};

% Conexões
\draw[seta] (rsl) -- (exp);
\draw[seta] (exp) -- (modelo);
\draw[seta] (rsl.south) |- ([yshift=0.3cm]decisao.west);
\draw[seta] (exp.south) -- (decisao.north);
\draw[seta] (modelo.south) |- ([yshift=0.3cm]decisao.east);

\end{tikzpicture}
\caption{Encadeamento conceitual da pesquisa.}
\label{fig:fluxo-pesquisa}
\end{figure}

\section{Objetivos}
\label{sec:intro-objetivos}

O objetivo geral desta pesquisa é investigar como evidências provenientes da literatura científica, resultados experimentais e características observáveis dos dados podem ser integrados para apoiar a tomada de decisão arquitetural em sistemas intensivos em dados.

Para alcançar esse objetivo geral, foram definidos os seguintes objetivos específicos:

\begin{enumerate}
\item caracterizar, por meio de uma Revisão Sistemática da Literatura, as principais técnicas, abordagens, métricas e lacunas relacionadas ao design de sistemas intensivos em dados;


\item sintetizar evidências quantitativas sobre o impacto de otimizações arquiteturais e sistêmicas no desempenho de sistemas intensivos em dados;

\item projetar e executar uma campanha experimental para comparar estratégias de processamento \textit{batch} e \textit{streaming} sob diferentes condições operacionais;

\item analisar estatisticamente os resultados experimentais, considerando métricas como \textit{throughput}, latência, escalabilidade, uso de recursos e estabilidade do processamento;

\item propor uma abordagem quantitativa de apoio à decisão arquitetural que relacione características dos dados e do processamento ao comportamento esperado do sistema.


\end{enumerate}

A justificativa desta pesquisa reside na necessidade de reduzir a dependência exclusiva de avaliações empíricas isoladas no processo de decisão arquitetural. Em sistemas intensivos em dados, a realização de experimentos completos para cada alternativa arquitetural pode demandar tempo, infraestrutura computacional e esforço de implementação significativos. Dessa forma, uma abordagem baseada em evidências pode auxiliar arquitetos de software na análise preliminar de alternativas, contribuindo para decisões mais fundamentadas, rastreáveis e alinhadas aos requisitos de desempenho e escalabilidade do sistema.

Além disso, esta pesquisa contribui para a área de Engenharia de Software ao aproximar três dimensões complementares: a consolidação de evidências científicas por meio de revisão sistemática, a produção de evidências empíricas por meio de experimentação controlada e a construção de mecanismos quantitativos para apoiar decisões arquiteturais. Essa integração busca avançar o conhecimento sobre o design de sistemas intensivos em dados e fornecer subsídios para pesquisas futuras relacionadas à avaliação, predição e recomendação de arquiteturas orientadas ao processamento intensivo de dados.


Esta proposta está organizada em oito capítulos. O Capítulo~\ref{cap:fundamentacao} apresenta a fundamentação teórica relacionada a sistemas intensivos em dados, decisões arquiteturais, processamento \textit{batch} e \textit{streaming}, Apache Spark, Apache Kafka e Teoria da Informação. O Capítulo~\ref{cap:metodologia} descreve a metodologia da pesquisa. O Capítulo~\ref{cap:rsl} apresenta a Revisão Sistemática da Literatura conduzida nesta pesquisa. O Capítulo~\ref{cap:campanha-experimental} descreve a campanha experimental realizada. O Capítulo~\ref{cap:resultados} apresenta os resultados parciais. O Capítulo~\ref{cap:cronograma} apresenta o cronograma da pesquisa. Por fim, o Capítulo~\ref{cap:proximos-passos} descreve os próximos passos e as atividades previstas para continuidade da dissertação.
